% -*- coding: utf-8 -*-
% !TEX program = xelatex
\documentclass[plain,noanswer,medmath]{../../template/jnuexam}
\usepackage{../../template/kysx}

\begin{document}


\examtitle{2009}{3}


\loadproblems{1}{2009P1}%载入试卷一
\loadproblems{2}{2009P2}%载入试卷二


\makepart{选择题}{1～8小题，每小题4分，共32分}


\useproblem{2}{1}{1}%【同试卷二第一(1)题】


\useproblem{1}{1}{1}%【同试卷一第一(1)题】


\begin{problem}
使不等式$\int_1^x \frac{\sin t}{t}\dt > \ln x$成立的$x$的范围是\pickout{}
\begin{abcd}
\item $(0,1)$．
\item $\left(1,\frac{\pi}{2} \right)$．
\item $\left(\frac{\pi}{2},\pi \right)$．
\item $(\pi , +\infty)$．
\end{abcd}
\end{problem}


\useproblem{1}{1}{3}%【同试卷一第一(3)题】


\useproblem{1}{1}{6}%【同试卷一第一(6)题】


\useproblem{2}{1}{8}%【同试卷二第一(8)题】


\begin{problem}
设事件$A$与事件$B$互不相容，则\pickout{}
\begin{abcd}
\item $P(\overline{A}\;\overline{B}) = 0$．
\item $P(AB) = P(A)P(B)$．
\item $P(A) = 1 - P(B)$．
\item $P(\overline{A} \cup \overline{B}) = 1$．
\end{abcd}
\end{problem}


\useproblem{1}{1}{8}%【同试卷一第一(8)题】


\makepart{填空题}{9～14小题，每小题4分，共24分}


\begin{problem}
$\lim_{x \to 0} \frac{\e - \e^{\cos x}}{\sqrt[3]{1 + x^2} - 1} =$ \fillin{}．
\end{problem}


\begin{problem}
设$z = (x + \e^y)^x$，则$\left. \frac{\pdz}{\pdx} \right|_{(1,0)} =$ \fillin{}．
\end{problem}


\begin{problem}
幂级数$\sum_{n = 1}^{\infty}\frac{\e^n - (- 1)^n}{n^2}x^n$的收敛半径为 \fillin{}．
\end{problem}


\begin{problem}
设某产品的需求函数为$Q = Q(p)$，其对价格$p$的弹性$\varepsilon_p = 0.2$，
则当需求量为$10000$件时，价格增加$1$元会使产品收益增加\fillin{}元．
\end{problem}


\begin{problem}
设$\alpha = (1,1,1)^T$,$\beta = (1,0,k)^T$．若矩阵$\alpha \beta^T$相似于$\begin{pmatrix}
  3 & 0 & 0 \\
  0 & 0 & 0 \\
  0 & 0 & 0
\end{pmatrix}$，则$k =$ \fillin{}．
\end{problem}


%【类似试卷一第二(14)题】
\begin{problem}
设$X_1,X_2, \cdots ,X_m$为来自二项分布总体$B(n,p)$的简单随机样本，
$\overline{X}$和$S^2$分别为样本均值和样本方差，
记统计量$T = \overline{X} - S^2$，则$ET =$ \fillin{}．
\end{problem}


\makepart{解答题}{15～23小题，共94分}


\useproblem[points=9]{1}{3}{15}%【同试卷一第三(15)题】


\useproblem[points=10]{2}{3}{16}%【同试卷二第三(16)题】


\useproblem[points=10]{2}{3}{19}%【同试卷二第三(19)题】


\useproblem[points=11]{1}{3}{18}%【同试卷一第三(18)题】


\begin{problem}[points=10]
设曲线$y = f(x)$，其中$f(x)$是可导函数，且$f(x) > 0$．
已知曲线$y = f(x)$与直线$y = 0,x = 1$及$x = t$（$t > 1$）所围成的曲边梯形绕$x$
轴旋转一周所得的立体体积值是该曲边梯形面积值的$\pi t$倍，求该曲线方程．
\end{problem}


\useproblem[points=11]{1}{3}{20}%【同试卷一第三(20)题】


\useproblem[points=11]{1}{3}{21}%【同试卷一第三(21)题】


\begin{problem}[points=11]
设二维随机变量$(X,Y)$的概率密度为
$$f(x,y) = \begin{cases}
  \e^{- x}, & 0 < y < x, \\
         0, & \text{其他}.
\end{cases}$$
\begin{enumerate*}
\item 求条件概率密度$f_{Y|X}(y|x)$；
\item 求条件概率$P\left\{X \le 1 \middle| Y \le 1 \right\}$．
\end{enumerate*}
\end{problem}


\useproblem[points=11]{1}{3}{22}%【同试卷一第三(22)题】


\end{document}
